\section{Introduction}\label{sec:intro}

How large is the fiscal multiplier, and what determines its size? A substantial
empirical and theoretical literature has identified three facts that any serious
model of fiscal transmission must confront simultaneously \citep{Auclert2024}:

\begin{description}[leftmargin=1.4em]
  \item[Fact 1.] \textit{Marginal propensities to consume (MPCs) are high.}
    Households spend approximately 25 cents of every additional dollar of income
    in the quarter they receive it, and around 50 cents annually
    \citep{KaplanViolante2014,Fagereng2021}.
  \item[Fact 2.] \textit{Marginal propensities to earn (MPEs) are low.} When
    households receive a windfall, they reduce their labor supply very little;
    empirical estimates place the annual MPE between $0$ and $0.04$
    \citep{Cesarini2017}.
  \item[Fact 3.] \textit{Fiscal multipliers are moderate.} Under accommodative
    monetary policy, both impact and cumulative multipliers lie in the range
    $0.6$--$2.0$ \citep{Ramey2019,Christiano2011,Woodford2011}.
\end{description}

\noindent\auclert{} show that these three facts create a \emph{trilemma} for New
Keynesian models with frictionless labor supply: it is impossible to match all
three simultaneously with standard separable preferences. Matching Facts~1 and~2
requires strong consumption--labor complementarity ($U_{cn}\neq 0$), but this
same complementarity generates multipliers that violate Fact~3. \BHL{} resolve
the trilemma through a different non-separability: by separately parameterizing
complementarity and income effects in $U(c,n)$, they allow the model to target
MPCs and MPEs independently of the multiplier.

\smallskip
This paper introduces a \emph{new dimension} to this debate. Both \auclert{} and
\BHL{} work with utility over consumption and labor: $U(c,n)$. We extend the
utility function to $U(c,g,n)$, where $g$ is government spending -- a publicly
provided good that directly enters household preferences. The cross-partial
$U_{cg}^{j}\neq 0$ is \emph{\EC{}} between private and public consumption: a
rise in $g$ raises (if $U_{cg}^{j}>0$) or lowers (if $U_{cg}^{j}<0$) the marginal
utility of private consumption. This is empirically well-documented
\citep{Karras1994,JallesKarras2022,FrancoisKeinsley2023,Francois2023,BarthelFrancois2025}
and generates a fiscal-transmission mechanism that is entirely absent from
models where government spending enters only through the resource constraint.

\paragraph{Contribution.} Our central contribution is threefold, with each point
corresponding to an empirically-testable new channel:

\begin{enumerate}[leftmargin=1.6em]
  \item We characterize a \textbf{direct fiscal channel $\xi$} -- absent from
    the separable framework of \bilbiie{} -- that captures how government
    spending directly stimulates private demand through its effect on the
    marginal utility of consumption. This channel depends on the
    cross-elasticities $\epsilon^{j}_{cg}$ and is new to the literature.
    Crucially, $\xi$ and the standard income-cyclicality statistic $\chi$ are
    cleanly separated: $\chi$ is identified by $\epsilon^{j}_{cc}$ moments alone,
    while $\xi$ is identified by $\epsilon^{j}_{cg}$ moments -- giving
    researchers two independent sufficient statistics for fiscal transmission.
  \item We show that \textbf{heterogeneous internalization}
    ($\theta^{H}\neq\theta^{S}$) generates curvature heterogeneity across
    household types. Since $\epsilon^{j}_{cc}=\sigma\alpha^{j}$ under the linear
    specification, $\theta^{H}\neq\theta^{S}$ induces
    $\alpha^{H}\neq\alpha^{S}$, which shifts income cyclicality $\chi$
    \emph{independently} of the hand-to-mouth share $\lambda$. This $\theta$-induced
    curvature heterogeneity provides a new route to the MPC/MPE trilemma
    resolution of \auclert{} -- one that operates through the government
    spending margin rather than through consumption--labor non-separability,
    and works even when $U_{cn}=0$ exactly.
  \item We derive a \textbf{distributional dominance condition} linking
    inequality, complementarity, and multiplier size: fiscal policy is most
    effective precisely where it is most needed -- in unequal economies where
    the poor disproportionately complement public goods with private
    consumption. Quantitatively, in a HANK calibration disciplined by US micro
    data, moving from the separable benchmark to heterogeneous internalization
    raises impact multipliers by 0.12--0.34 log points.
\end{enumerate}

\paragraph{The direct fiscal channel: a new sufficient statistic.}
The first contribution deserves emphasis because it stands independently of the
trilemma discussion and is arguably the paper's cleanest result. In \bilbiie{}'s
separable framework, the fiscal multiplier depends on one sufficient statistic:
income cyclicality $\chi$. Government spending affects output only through
general-equilibrium income effects -- it raises wages and profits, which raise
\HtM{} income, which feeds back into demand. There is no channel by which $g$
directly affects the private consumption decision.

\smallskip
Introducing $U_{cg}^{j}\neq 0$ creates a second sufficient statistic: $\xi$, the
direct fiscal channel. When $g$ rises, the marginal utility of private
consumption changes -- upward for households with $\epsilon^{j}_{cg}<0$
(complements) and downward for those with $\epsilon^{j}_{cg}>0$ (substitutes).
This directly alters how much \HtM{} households spend out of any given income,
above and beyond the standard wealth and income effects. Formally, $\xi$ captures
the direct response of \HtM{} consumption to the fiscal shock:
$\hat c^{H}_{t}=\chi\,\hat c_{t}+\xi\,\hat g_{t}$, where the $\xi\,\hat g_{t}$
term is entirely new. In \bilbiie{}'s separable world,
$\xi=\zeta(\chi-\mu/\lambda)$ -- a function of $\chi$ alone. In our framework,
$\xi$ also depends on $\epsilon^{H}_{cg}$ and $\epsilon^{S}_{cg}$, giving it
\emph{independent variation}.

\smallskip
Crucially, $\xi$ and $\chi$ are cleanly separated in the general specification:
Proposition~\ref{prop:chi-general} shows $\chi$ depends on $\epsilon^{j}_{cc}$
only, while Proposition~\ref{prop:xi-general} shows $\xi$ depends on
$\epsilon^{j}_{cg}$. This separation means the two channels can be
\emph{independently disciplined by data} -- income cyclicality moments identify
$\epsilon^{j}_{cc}$, while the consumption response to government spending
shocks identifies $\epsilon^{j}_{cg}$. The distributional dominance condition
in Corollary~\ref{cor:het} follows directly from the $\xi$ channel: whether
fiscal policy generates larger or smaller multipliers than the separable
benchmark depends on whether the \HtM{} complementarity effect (which raises
$\xi$ when $\epsilon^{H}_{cg}<0$) dominates the saver substitutability effect
(which dampens $\xi$ when $\epsilon^{S}_{cg}>0$). This is a result that has
nothing to do with the trilemma and holds regardless of whether
$\theta^{H}=\theta^{S}$.

\paragraph{The new trilemma resolution.}
In the standard separable THANK framework \citep{Bilbiie2025}, income cyclicality
is $\chi^{\mathrm{sep}}=1+\phi(1-\tau_{D}/\lambda)$, which depends on $\lambda$
(the \HtM{} share). Since the aggregate MPC also depends primarily on $\lambda$,
the two are coupled: raising $\lambda$ to match high MPCs simultaneously raises
$\chi$ (MPEs). This is the separable version of the trilemma within THANK.

\smallskip
Now introduce $\theta^{H}\neq\theta^{S}$ -- heterogeneous internalization. Under
Assumption~\ref{asm:linear}, $\epsilon^{j}_{cc}=\sigma\alpha^{j}$ where
$\alpha^{j}=\bar c/(\bar c+\theta^{j}\bar g)$. When $\theta^{H}\neq\theta^{S}$,
this induces $\alpha^{H}\neq\alpha^{S}$ -- curvature heterogeneity across
household types. Income cyclicality becomes
\[
  \chi(\theta^{S},\theta^{H})=
  \frac{K(\phi+\sigma\alpha^{S})}
       {(\phi+\sigma\alpha^{H}) - K\sigma\lambda(\alpha^{H}-\alpha^{S})}.
\]
The aggregate MPC is still governed primarily by $\lambda$. But $\chi$ now also
depends on $\alpha^{H}-\alpha^{S}$ -- the $\theta$-induced curvature gap --
\emph{independently} of $\lambda$. By adjusting $\theta^{H}-\theta^{S}$ while
holding $\lambda$ fixed, a researcher can shift $\chi$ to match Fact~2 without
disturbing the MPC calibration for Fact~1.

\smallskip
Two clarifications are essential. First, the mechanism runs through
$\epsilon^{j}_{cc}$, not through $\epsilon^{j}_{cg}$ directly: $\chi$ does not
depend on $\epsilon^{j}_{cg}$. The resolution is therefore more precisely a
$\theta$-induced curvature-heterogeneity mechanism -- $\theta^{j}$ shifts the
effective consumption curvature $\sigma\alpha^{j}$, and it is this curvature gap
that decouples $\chi$ from $\lambda$. Second, this works even with fully
separable consumption--labor preferences ($U_{cn}=0$), making it
\emph{orthogonal} to the resolutions of \auclert{} and \BHL{}.

\paragraph{Why complementarity is heterogeneous.}
The empirical evidence strongly motivates $\theta^{H}\neq\theta^{S}$ as a
maintained hypothesis. \citet{Gethin2023} constructs novel estimates of global
poverty that incorporate the consumption of public services and documents that
public goods are strongly progressive in their incidence: poor households are
disproportionate beneficiaries of public education, healthcare, housing, and
transport infrastructure. Public goods represent 30\% of global GDP and have
accounted for approximately 20\% of global poverty reduction since 1980. This
distributional incidence is directly consistent with \HtM{} households treating
public goods as complements to private consumption: public health services raise
the productivity of private spending on nutrition; public transport enables
private labor-market participation and reduces commuting costs for low-income
workers; publicly provided education complements private inputs into children's
human capital in low-income households. For these agents, a rise in $g$ raises
$U^{H}_{c}$, inducing more private spending on top of the direct demand effect.
\citet{Francois2023} finds structural evidence for this complementarity among
rule-of-thumb consumers in Sub-Saharan Africa; \citet{BarthelFrancois2025}
establish that the complementarity relationship depends on income through
non-homothetic preferences, motivating type-varying $\theta^{j}$.

Wealthier unconstrained households, by contrast, tend to substitute: they
privately purchase alternatives -- private healthcare, private schools, private
security, private transport -- so public expansion may crowd out their private
spending. \citet{Gethin2023} documents this aggregate substitution pattern
across countries: where public provision is lower, private out-of-pocket spending
on healthcare and education is higher. This asymmetry is precisely what
generates the curvature gap: $\theta^{H}<0$ (complements) raises $\alpha^{H}$
while $\theta^{S}>0$ (substitutes) lowers $\alpha^{S}$, widening
$\alpha^{H}-\alpha^{S}$ and shifting $\chi$ independently of $\lambda$.

\paragraph{Empirical discipline.}
Section~\ref{sec:empirical} develops a micro-to-macro identification strategy
that delivers household-type-specific estimates of $\theta^{j}$ suitable for
structural calibration. We combine three sources of variation. First, we use
household-level cross-elasticities between public-good provision and private
consumption estimated from the \textit{Consumer Expenditure Survey} (CEX) matched
to state-year variation in Medicaid expansions, public-transit funding, and
K--12 spending. Second, we exploit the quasi-experimental variation in the
incidence of public goods documented by \citet{Gethin2023} to discipline the
aggregate dispersion across household types. Third, we use the VAR-based
consumption response to government-spending shocks \citep{Ramey2011} as a
pattern-matching target for the reduced-form $\xi$ statistic. Our preferred
estimates imply $\theta^{H}\in[-0.90,-0.62]$ and
$\theta^{S}\in[0.21,0.38]$ -- a curvature gap that is statistically and
economically significant.

\paragraph{Main analytical results.}
We embed $U(c,g,n)$ with $U_{cg}^{j}\neq 0$ into the THANK framework of
\bilbiie{} and derive closed-form results under general non-separable
preferences, then specialize to the linear effective consumption form
$\tilde c^{j}=c^{j}+\theta^{j}g$. Two sufficient statistics characterize the
equilibrium. Income cyclicality $\chi$ (the MPE in our model) is determined by
own-consumption curvature $\epsilon^{j}_{cc}$ and by the heterogeneous
internalization gap $\theta^{H}-\theta^{S}$. The direct fiscal channel $\xi$ is
determined by the cross-elasticities $\epsilon^{j}_{cg}$ capturing \EC{}. These
two objects are cleanly separated in the symmetric case: $c$--$g$
non-separability modifies $\xi$ but leaves $\chi$ equal to \bilbiie{}'s
expression. Under heterogeneous internalization, the separation breaks and
$\chi$ acquires additional degrees of freedom.

The general fiscal multiplier is
\begin{equation}\label{eq:multiplier-intro}
  \mu(\theta^{S},\theta^{H}) = 1 +
  \frac{\lambda\,\xi(\theta^{S},\theta^{H})}
       {1-\lambda\,\chi(\theta^{S},\theta^{H})}.
\end{equation}
A distributional dominance condition governs the net effect: fiscal policy is
more effective in more unequal economies \emph{both} because of high MPCs among
the poor and because the poor are precisely those for whom $g$ is most
complementary to private consumption. Countercyclical inequality
($\chi>1$, empirically standard: \citealp{Bilbiie2025,Patterson2023}) amplifies
both channels simultaneously.

\paragraph{Quantitative results.}
Section~\ref{sec:quant} takes the model to a full HANK calibration following
\citet{KaplanMollViolante2018}. We discipline the wealth distribution to match
the Survey of Consumer Finances, the intertemporal MPC to match
\citet{Fagereng2021}, and the distribution of $\theta^{j}$ using the empirical
estimates from Section~\ref{sec:empirical}. Three findings stand out. First, the
impact multiplier rises from $1.50$ in the separable benchmark to $1.68$--$1.84$
under heterogeneous internalization -- economically large relative to the
dispersion in the VAR-based empirical literature. Second, the multiplier
amplification is concentrated in the bottom of the wealth distribution: the
consumption response in the bottom quintile is 2.3 times larger than in the top
quintile, versus a ratio of 1.4 in the separable model. Third, scenario analysis
reveals that fiscal policy targeted at public services used by low-income
households (health, transit, education) generates output multipliers up to 38\%
larger than untargeted transfers.

\paragraph{Policy implications.}
The model transforms fiscal policy from a question of aggregate size into a
question of distributional design. The stakes are large: \citet{Gethin2023}
documents that public goods account for 20\% of global poverty reduction since
1980. Our framework provides a general-equilibrium mechanism that helps explain
why: spending that reaches complementarity-experiencing households generates
larger aggregate demand effects, aligning equity and efficiency. Fiscal
consolidation has \emph{asymmetric} welfare costs -- cutting public services
depresses $U^{H}_{c}$ for constrained households, compressing their effective
demand in ways that standard models systematically underestimate, because those
models suppress the $c$--$g$ complementarity channel entirely. The welfare cost
of austerity falling on public health, education, and transport is therefore
larger than conventional multiplier estimates suggest.

\smallskip
This is particularly acute in low-income and developing economies.
\citet{Gethin2023} documents a ``triple curse'' for poor countries: they spend
less on public services, spend more regressively, and provide lower quality.
In our framework, these are precisely the economies where $\lambda$ is largest,
where $|\theta^{H}|$ is potentially strongest (poor households rely most heavily
on public services as inputs to private consumption), and where the curvature
gap $\alpha^{H}-\alpha^{S}$ is widest. The model therefore predicts that
well-targeted fiscal expansions in low-income countries -- directing public
spending toward services used by constrained households -- could generate
multiplier effects substantially larger than those estimated for advanced
economies, precisely because the distributional dominance condition is most
favorable there. Automatic stabilizers that maintain public consumption during
downturns do double duty: they stabilize aggregate demand and protect the
private consumption complementarity of the most vulnerable, a consideration
entirely absent from models where $g$ enters only through the resource
constraint.

\paragraph{Related literature.}
Our paper contributes to five strands. \emph{First} and most directly, the
literature on the MPC/MPE trilemma: \citet{Auclert2024} who identify it and
propose sticky wages as the solution; \BHL{} who resolve it through $c$--$n$
complementarity ($U_{cn}\neq 0$); we provide a complementary resolution through
$\theta$-induced curvature heterogeneity under heterogeneous internalization
($\theta^{H}\neq\theta^{S}$), which operates through the government-spending
margin and is orthogonal to existing labor-supply-based resolutions.
\emph{Second}, the literature on non-separable government spending in utility:
\citet{BouakezRebei2007}, \citet{LinnemannSchabert2006}, \citet{Feve2013} in
RANK; \citet{Bilbiie2011} on the fiscal policy puzzle; and the empirical
literature on public--private complementarity
\citep{Karras1994,JallesKarras2022,FrancoisKeinsley2023,FrancoisDawood2018,
Francois2023,BarthelFrancois2025}. \emph{Third}, the HANK literature on fiscal
multipliers \citep{KaplanMollViolante2018,Bilbiie2025,AuclertIKC2024,Broer2023}.
\emph{Fourth}, heterogeneous-agent fiscal transmission
\citep{Gali2007,OhReis2012,Hagedorn2019}. \emph{Fifth}, the distributional
incidence of public goods and its macroeconomic implications \citep{Gethin2023},
which provides the empirical grounding for the heterogeneous-internalization
assumption. To the best of our knowledge, this is the first paper to: (i) derive
closed-form fiscal multipliers under $c$--$g$ \EC{} in THANK; (ii) show that
$\theta$-induced curvature heterogeneity under heterogeneous internalization
provides a new route to the MPC/MPE trilemma resolution through the
government-spending margin; and (iii) characterize a distributional dominance
condition linking inequality, complementarity, and the size of the fiscal
multiplier, with direct implications for fiscal policy in unequal and
developing economies.

\paragraph{Roadmap.}
Section~\ref{sec:model} presents the model. Section~\ref{sec:analytical} derives
the analytical results. Section~\ref{sec:quant} presents the quantitative
analysis. Section~\ref{sec:empirical} develops the empirical identification
strategy. Section~\ref{sec:policy} conducts policy experiments and
counterfactuals. Section~\ref{sec:conclusion} concludes.
Appendices~\ref{app:fullderivation}--\ref{app:calibration} provide full
derivations, alternative proofs, the complete model setup, and calibration and
simulation code.
